
\chapter{Akt II -- Das Auge der Wüste}
\section{Szene 1 -- Aufbruch ins Ungewisse}
\begin{meisterinfo}
\textbf{Szenenziel.} Die Ermittler verlassen das vertraute Arkham, lernen den Alltag einer langen Reise von 1920 kennen und begegnen auf dem Atlantik Yusuf Hassan. Er ist ein vermisster Arbeiter aus Flynns Expedition, aber seine Erscheinung ist keine gewöhnliche Reisebekanntschaft. Die Szene soll neugierig machen, nicht den Mythos bereits erklären.

\textbf{Historischer Anker.} Die SS \emph{Adriatic} der White Star Line ist mit einer Abfahrt von New York nach Cherbourg und Southampton am \textbf{24. April 1920} in erhaltenen Passagierlisten belegt. Alle übrigen genauen Tage sowie die Anschlüsse in Europa sind eine \emph{plausible Szenario-Festlegung}, kein historisch belegter durchgehender Reiseplan. Alternativ kann die Spielleitung Schiff und Datum frei halten.

\textbf{Vorgeschlagene Chronologie:} 5. April Treffen mit Garmond; 6.--22. April Untersuchungen, Ausrüstung, Visa, Schiffsbuchung und Anreisevorbereitung; 23. April Bahn nach New York; 24. April Abfahrt der \emph{Adriatic}; Ende April/Anfang Mai Ankunft in Europa; Anfang Mai Bahn- und Kanalreise nach Marseille; darauf etwa fünf bis sieben Tage Mittelmeerfahrt mit einem passenden Passagierdampfer; um die zweite Maiwoche Alexandria und Weiterreise nach Kairo. Die genauen Anschlussdaten bleiben flexibel.

\textbf{Quellen zur historischen Einordnung:} GG Archives, \emph{White Star Line Passenger Lists 1889--1939}, Eintrag \emph{SS Adriatic, 24 April 1920}; \emph{Cairo of To-day}, Reisehinweise zur Seeverbindung Marseille--Alexandria (1897/98; Beleg der historischen Route, nicht eines Fahrplans für 1920).
\end{meisterinfo}

\subsection{1.1 -- Abschied von Arkham}
\begin{vorlesetextalle}
Die Tage nach Eurem Besuch in der Universitätsbibliothek vergehen zwischen Vorbereitungen, Briefwechseln und geschäftigen Besorgungen. Die Miskatonic University übernimmt die Kosten der Reise. Koffer werden gepackt, Unterlagen geordnet, dringend benötigte Ausrüstung beschafft. Zwischen den vielen kleinen Erledigungen bleibt immer wieder Zeit, an das Haus der Carlsons zurückzudenken: an den seltsam staubfreien Tisch, die Frau unter dem Dach und die Kreatur, die Euch aus einer Vergangenheit heraus anzusehen schien.

Als Ihr schlie\ss{}lich an einem kühlen Apriltag den Zug Richtung New York besteigt, liegt Arkham hinter Euch. Fabrikschornsteine, kleine Bahnhöfe und graue Vorstadtstra\ss{}en ziehen am Fenster vorbei. Vor Euch liegen ein Ozean, ein fremder Kontinent und die Suche nach Männern, deren letzte Nachrichten alles andere als beruhigend waren.
\end{vorlesetextalle}
\begin{meisterinfo}
\textbf{Spielraum:} Gib jeder Figur die Gelegenheit, eine persönliche Vorbereitung zu schildern. Nützliche Vorbereitungen sind Reisepapiere, Kleidung für die Wüste, Notizmaterial, Medikamente und eine sichere Aufbewahrung für Briefe und Karte. Gewöhnliche Ausrüstung soll nicht an Würfen scheitern. Eine \textbf{Bibliotheksnutzung}- oder \textbf{Wissen}-Probe darf Informationen zu den Verhältnissen in Ägypten liefern, aber keine zentrale Lösung vorwegnehmen.

\textbf{Kosten:} Gemä\ss{} Garmonds Auftrag finanziert die Universität Anreise, Passage und übliche Ausrüstung. Die Höhe etwaiger persönlicher Ausgaben ist für diese Szene nicht entscheidend.
\end{meisterinfo}

\subsection{1.2 -- Die Atlantiküberquerung}
\begin{vorlesetextalle}
New York empfängt Euch mit Lärm, Autohupen, den Rufen der Hafenarbeiter und dem unaufhörlichen Pfeifen der Dampfschiffe. Zwischen Kr\"anen und Lagerhäusern erhebt sich Euer Ozeandampfer: ein Schiff aus Stahl, mit hohen Schornsteinen, Rettungsbooten und Decks, die scheinbar nicht enden wollen.

Am 24. April verlasst Ihr den Hafen. Das Ufer fällt langsam zurück, bis die Stadt nur noch eine zackige Linie am Horizont ist. Die ersten Tage an Bord haben einen beinahe beruhigenden Rhythmus. Mahlzeiten werden zu festen Stunden serviert. Im Salon werden Zeitungen gelesen, an Deck promenieren Passagiere unter Decken und Mänteln, und das dumpfe Schlagen der Maschinen begleitet Euch bis in den Schlaf.

Im Lauf der Überfahrt verschwinden die letzten sichtbaren Spuren des Festlandes. Ringsum liegt nur der graublaue Atlantik. Selbst die Gespräche an Bord scheinen manchmal leiser zu werden, wenn die See besonders still ist.
\end{vorlesetextalle}
\begin{meisterinfo}
\textbf{Atmosphäre und Interaktion:} Kurze Gespräche mit anderen Reisenden, ein Rundgang über das Deck, eine Partie Karten oder ein Blick in die Schiffsbibliothek sind möglich. Erzähle einige Tage in geraffter Form; mache daraus keine zweite gro\ss{}e Ermittlungsphase. Auf dem Dampfer gibt es keine festgelegte Mythosbedrohung au\ss{}er der folgenden Begegnung.

\textbf{Ansto\ss{} zur Nebelszene:} Am sechsten Abend der Fahrt können die Ermittler nach dem Essen noch einmal auf das Promenadendeck gehen. Niemand muss eine Probe schaffen, um die Begegnung zu erleben. Falls alle in ihren Kabinen bleiben, bittet ein Steward um Hilfe bei einer offenstehenden Deckstür und lockt sie so nach drau\ss{}en, ohne die Situation zu erzwingen.
\end{meisterinfo}

\subsection{1.3 -- Der Fremde im Nebel}
\begin{vorlesetextalle}
Es ist die sechste Nacht Eurer Atlantiküberquerung. Seit dem Nachmittag liegt dichter Nebel über dem Meer. Anfangs waren es nur wei\ss{}e Schwaden gewesen, die über das Wasser zogen. Nun hat sich eine undurchdringliche graue Wand um das Schiff gelegt.

Das regelmä\ss{}ige Dröhnen der Maschinen reicht bis auf das Promenadendeck. In grö\ss{}eren Abständen ertönt das tiefe Nebelhorn. Der Klang verliert sich in der Dunkelheit, ohne ein Echo zurückzuwerfen.

Die meisten Passagiere haben sich längst in die Innenräume zurückgezogen. Die feuchten Decksplanken glänzen im Licht der vereinzelten Lampen. Am äu\ss{}eren Geländer steht ein Mann, den Ihr hier zuvor noch nicht bemerkt habt.

Er trägt einfache, abgetragene Kleidung und hat ein dunkles Tuch um seine Schultern geschlungen. Sein Rücken ist Euch zugewandt. Er starrt hinaus in den Nebel, als versuche er darin etwas zu erkennen.

Als Ihr näher kommt, dreht er langsam den Kopf. Sein Gesicht ist schmal und erschöpft, die Augen liegen tief in ihren Höhlen. Er mustert Euch einen Augenblick schweigend. Dann sagt er mit rauer, leiser Stimme:

\glqq{}Ihr sucht den Mann mit dem Stein.\grqq{}

Es klingt nicht wie eine Frage.

\glqq{}Er glaubt, er hätte etwas gefunden. Aber er hat etwas geöffnet.\grqq{}

Der Mann schaut hastig über die Schulter, als hätte jemand hinter ihm gesprochen.

\glqq{}Sagt Carlson, dass er den Augen nicht trauen darf.\grqq{}

Er tritt einen Schritt zurück. Der Nebel zieht in dichten Fäden zwischen Euch und das Geländer.

\glqq{}Mein Name ist Yusuf. Yusuf Hassan. Vergesst ihn nicht.\grqq{}

Ein langer Sto\ss{} des Nebelhorns fährt Euch durch Mark und Bein. Für einen Moment verschluckt der wirbelnde Nebel die Gestalt. Als Ihr wieder an dieselbe Stelle blickt, ist dort niemand mehr.
\end{vorlesetextalle}
\begin{meisterinfo}
\textbf{Wer ist Yusuf?} Yusuf Hassan ist ein ägyptischer Arbeiter aus der ursprünglichen Expedition. Flynn erwähnt in seinem Brief vom 7. Februar 1920 einen Arbeiter, der einem Schrei in die Anlage folgte und verschwand. \emph{Für unsere neue Fassung} erhält dieser bislang namenlose Mann den Namen Yusuf Hassan. Sein Tod ist nicht bestätigt.

\textbf{Spielweise:} Halte für Rückfragen einige knappe Sätze bereit. Auf \glqq{}Woher kennen Sie uns?\grqq{} erwidert er: \glqq{}Ich kenne nicht Euch. Aber der Stein schaut.\grqq{} Auf \glqq{}Wo ist Flynn?\grqq{} antwortet er: \glqq{}Hinter den Türen, die keine Türen sind.\grqq{} Bei Berührung ist sein Ärmel feucht und eiskalt. Die Begegnung endet, wenn sich der Nebel schlie\ss{}t oder ein Schiffssignal erklingt. Es gibt keine Verfolgungsjagd über die Reling; Yusuf verschwindet zwischen zwei Herzschlägen.

\textbf{Stabilität:} Noch keine Probe beim blo\ss{}en Gespräch. Der spätere Nachweis, dass Yusuf nicht an Bord war, rechtfertigt \textbf{0/1W3 STA}. Aufmerksame Ermittler können merken, dass er keine Schuhe trägt und trockene Sandkörner an den Fü\ss{}en haften. Diese Beobachtung erfordert \textbf{Verborgenes Erkennen (regulär)}, ist aber nicht nötig, um die Szene zu verstehen.
\end{meisterinfo}

\subsection{1.4 -- Ein Passagier, den es nicht gibt}
\begin{vorlesetextalle}
Einige Minuten später fährt das Schiff noch immer durch dichten Nebel. Ein Offizier der Nachtwache hört sich Euren Bericht mit zunehmendem Stirnrunzeln an. Er fragt nach Kleidung, Erscheinungsbild und Namen des Fremden.

\glqq{}Yusuf Hassan? Diesen Namen habe ich nie gehört. Warten Sie hier. Ich lasse nachsehen.\grqq{}

Nach einer Weile kehrt er mit einer Liste zurück. Er legt den Finger auf die eng beschriebenen Seiten und schüttelt langsam den Kopf.

\glqq{}Weder unter unseren Passagieren noch bei der Besatzung. Und ich versichere Ihnen: Während dieser Überfahrt ist niemand zugestiegen.\grqq{}
\end{vorlesetextalle}
\begin{meisterinfo}
\textbf{Nachforschungen:} Ein Gespräch mit dem Offizier und eine Prüfung des Verzeichnisses reichen aus, um Yusufs Abwesenheit zu bestätigen. \textbf{Überreden} oder \textbf{Charme} kann Zugang zum Originalverzeichnis verschaffen; auch ohne Erfolg wird die wesentliche Auskunft erteilt. \textbf{Verborgenes Erkennen (regulär)} an der Stelle der Begegnung entdeckt einige helle, ungewöhnlich trockene Sandkörner in den Fugen der nassen Planken. \textbf{Naturwissenschaft (Geologie)} kann vermuten, dass der feine Staub nicht typisch für eine salzige, feuchte Schiffsoberfläche ist; eine sichere geografische Herkunft lässt sich daraus nicht ableiten.

\textbf{Stabilität:} 0/1W3 STA, sobald die Ermittler begreifen, dass der Fremde dort nicht hätte sein können. Bei Fehlschlag auf soziale Proben wirkt der Offizier skeptischer, nicht unkooperativ. \textbf{Ausfallsicherung:} Wird der Name nicht notiert, lässt sich die Erinnerung später durch Aziz' Aussagen erneut auslösen.

\textbf{Spielleitergeheimnis:} Yusuf erscheint nicht als gewöhnlicher Geist. Er ist ein Echo, das durch die Störung von Raum und Zeit um den sehenden Stein entstanden ist. Damit wird auch die unmögliche Karte in der Bibliothek thematisch gespiegelt. Seine genaue Lage in den Katakomben bleibt für spätere Szenen offen; die Ermittler können ihn dort eventuell lebend, tot oder in einem Zwischenzustand finden. Das legen wir erst bei der Ausarbeitung der Grabkammer endgültig fest.
\end{meisterinfo}

\subsection{1.5 -- Über Europa}
\begin{vorlesetextalle}
Die letzten Tage auf dem Atlantik vergehen langsamer als zuvor. Immer wieder zieht es Euren Blick an jene Stelle des Decks, an der Yusuf gestanden hat. Doch weder im Licht der Morgensonne noch in weiteren Nebeln taucht er erneut auf.

Schlie\ss{}lich erscheint Land am Horizont. Europäische Hafenanlagen, Eisenbahnzüge, Zollbeamte und überfüllte Bahnsteige ersetzen die gleichförmige Weite des Ozeans. Ihr reist weiter Richtung Süden. Felder, Dörfer, Industriegebiete und gro\ss{}e Städte wechseln vor den Fenstern des Zuges, während sich mit jeder Etappe die Sprachen und die Gesichter um Euch verändern.

In Marseille liegt bereits das nächste Schiff vor Anker. Der salzige Geruch des Mittelmeeres, fremde Stimmen und die grelle Sonne über den Kais lassen die Atlantikreise wie einen weit entfernten Traum erscheinen.
\end{vorlesetextalle}
\begin{meisterinfo}
\textbf{Plausible Route:} Ankunft in Cherbourg beziehungsweise Southampton, dann Kanalquerung und Eisenbahnreise über Frankreich nach Marseille. Die konkrete Zugfolge und die Mittelmeerabfahrt sind bewusst nicht als historisch belegte April-/Mai-1920-Fahrpläne ausgegeben. Falls die Gruppe Europastopps ausspielen möchte, eignet sich eine kurze Begegnung mit einem Zollbeamten oder Reisevertreter, jedoch ohne neues Rätsel.
\end{meisterinfo}

\subsection{1.6 -- Das Mittelmeer und die Küste Ägyptens}
\begin{vorlesetextalle}
Das Mittelmeer zeigt sich Euch in einem anderen Licht. Das Wasser liegt heller und blauer unter einem beinahe makellosen Himmel. In der Ferne ziehen Frachter und Segelschiffe vorbei, über den Decks kreisen Möwen. Die Abende sind milder, die Luft riecht nach Salz und Maschinenöl.

Nach weiteren Tagen auf See erscheint schlie\ss{}lich die flache Küste Ägyptens. Zuerst sind nur helle Linien im Dunst zu erkennen, dann Kräne, Masten und die Häuser Alexandrias. Zwischen Rufen, Gepäckwagen und Zollformalitäten betretet Ihr zum ersten Mal ägyptischen Boden.

Die letzte Etappe führt mit der Eisenbahn nach Kairo. Durch das Fenster ziehen Felder, Dörfer und Wasserläufe vorbei. Hier und da öffnen sich Blicke auf eine Landschaft, deren Farben Ihr aus keinem Neuengland-Frühling kennt.

Als der Zug in den Bahnhof von Kairo einrollt, ist aus der Reise ein neues Kapitel geworden. Vor Euch liegt die Stadt, in der Flynn Carlson und Dr. Skorbinsk ihre Unternehmung begonnen haben.
\end{vorlesetextalle}
\begin{meisterinfo}
\textbf{Anschluss:} Akt II, Szene 2 -- \emph{Kairo, das Tor zum Orient}. Die Ermittler suchen zunächst das Hotel International auf. Dort soll der Kontakt zu Aziz wiederhergestellt werden.

\textbf{Spätere Rückbindung:} Wenn die Gruppe Aziz den Namen Yusuf Hassan nennt, erkennt er ihn als einen der vermissten Arbeiter wieder. Besonders wirkungsvoll ist ein altes Mannschaftsfoto, auf dem Yusuf zu sehen ist. Foto und Aussehen müssen erst in der Kairo-Szene festgelegt werden; hier ist noch kein zusätzliches Handout nötig.

\textbf{Leitmotiv:} Von der Karte \glqq{}Flynn, 1920\grqq{} über die Visionen im Carlson-Haus bis zu Yusufs Erscheinung auf See häufen sich Hinweise darauf, dass Orte und Zeiten rund um die Grabkammer nicht zuverlässig voneinander getrennt sind. Die Ursache wird in dieser Szene nicht erklärt.
\end{meisterinfo}
